\section{Introduction}\label{sec:intro}

Diabetic retinopathy (DR) affects about one in five people with diabetes, an estimated
103 million adults in 2020 and a projected 161 million by 2045\cite{teo2021prevalence}.
Deep learning systems that grade fundus photographs reached specialist-level agreement in
development studies\cite{gulshan2016,ting2017}, and autonomous screening devices have been
authorised on the strength of prospective pivotal trials\cite{abramoff2018pivotal,ipp2021eyeart}.
The same class of system behaves very differently once it leaves the population it was built
on: seven commercial algorithms applied to the same real-world images ranged from 51\% to 86\%
in referable-disease sensitivity\cite{lee2021headtohead}, and three algorithms evaluated against
one reference standard in Indian public-health clinics ranged from 14\% to 96\% in
specificity\cite{duggal2025realworld}.

Such failures are usually read as statements about what the network has learned. A recent
frozen-encoder benchmark found that two of three encoders, including the retinal foundation
model RETFound, did not identify a single mild case in an external corpus (grade-1 F1 of 0.000)
and attributed the result to the encoders' representations\cite{poyrazer2026calibration}.
An ordinal grading model does not output a grade, however. It outputs scores, and a
\emph{decision rule} --- a set of cut-points --- turns them into a grade and a referral. The
common default, 0.5 on every cumulative output or an argmax, is a choice that is rarely
reported or examined. If a model still ranks images correctly but its cut-points sit in the
wrong place for a new population, it will look blind to a grade it in fact separates, and two
models that differ mainly in where their defaults fall will look as if they differ in what they
learned.

Parts of this idea are not new. Recalibration before deployment is standard advice for
clinical prediction models\cite{janssen2008updating,vergouwe2017updating}. ORDER-DR, an ordinal
DR model developed on APTOS, explicitly distinguished threshold-dependent grade assignment from
threshold-independent risk ranking and noted that models trained on one public dataset can show
threshold shift externally\cite{sheng2026orderdr}. For binary referable DR across fundus
cameras, a recent preprint showed that discrimination transfers while probability calibration
degrades, mostly through prevalence, and that 100--200 labelled target images restore it
\cite{tabibu2026discrimination}; a public model repository reports that calibration degrades
unevenly across ordinal thresholds under domain shift\cite{gradeye2026}. What these studies do not
provide is a decomposition across many models and several populations of how much of an external
failure, and of an apparent difference between models, the ordinal decision rule accounts for;
evidence on why the score-to-grade mapping moves; a rule for comparing models that removes it;
and the label cost and risk of repairing it.

We provide these in a study whose analysis plans were committed to version control before the
corresponding analyses were run, reporting falsified predictions with the same weight as
confirmed ones. Ten development models trained on \pn{n.dev} images were evaluated on four
external populations --- APTOS 2019 (India), DDR (China, 147 hospitals), DeepDRiD (Shanghai)
and SUSTech-SYSU (Shenzhen) --- and a model trained on APTOS was evaluated in the reverse
direction. Two further models were fine-tuned on \pn{eye.eb3.trainn} EyePACS images instead of the
\pn{n.dev}-image development pool, to test whether the findings depend on the development set and
its labels; five seeds of the reference model measured training noise, and the released
predictions of ORDER-DR were reanalysed. The contributions are:
\begin{itemize}
\item a decomposition of external performance into discrimination at every ordinal threshold
and the mapping from score to grade, which shows that the mapping moves in opposite directions
on different populations through a shift in the score \emph{given} the true grade;
\item an account of what changing the decision rule can and cannot recover, including a negative
result for label-free prior adaptation;
\item the \emph{matched decision rule}, a comparison rule under which both models' cut-points are
refitted by the same procedure, with its record of changed, unchanged and reversed comparisons;
\item the number of locally labelled referable cases needed to set a referral threshold, with a
distribution-free probability of harm that translates into images at any screening prevalence, and
what the operating point means per 10\,000 people screened.
\end{itemize}
